\chapter{Intuitions on Cofinality}

\section{Distribution of Regular and Singular Cardinals}

To facilitate notation, we will say that a cardinal $\aleph_\alpha$ is a \textcolor{laser_eyes!80!black}{successor cardinal} if $\alpha$ is a successor ordinal, and we will say that it's a \textcolor{laser_eyes!80!black}{limit cardinal} if $\alpha$ is a limit ordinal.

Let's explore the concept we create with some cardinals. For the cases used in the motivation, we have that $\cf(\aleph_0)=\aleph_0$ and $\cf(\aleph_1)=\aleph_1$, i.e., both are regular cardinals, since they can't be written as a union of $\aleph_0$ or $\aleph_1$ smaller elements, respectively. In particular, the following lemma states that every successor cardinal $\aleph_{\alpha+1}$ is a regular cardinal, that is, it's not easy to "reach them from below", only with smaller sets.

\begin{lemma}{}{}
    Every successor cardinal $\aleph_{\alpha+1}$ is a regular cardinal.
\end{lemma}

\begin{proof}
    Assume by contradiction that $\aleph_{\alpha+1}$ is singular, by Lemma \ref{sing-sum}, it can be written as a sum of a smaller number of smaller cardinals:
    $$\aleph_{\alpha+1}=\sum_{i\in I}\kappa_i$$
    with $|I|<\aleph_{\alpha+1}$ and $\kappa_i<\aleph_{\alpha+1}$, $\forall i\in I$. Therefore $|I|,\kappa_i\leq\aleph_\alpha$, $\forall i\in I$, so, we have
    $$\aleph_{\alpha+1}=\sum_{i\in I}\kappa_i\leq\sum_{i\in I}\aleph_\alpha=\aleph_{\alpha}\cdot|I|\leq\aleph_\alpha\cdot\aleph_\alpha=\aleph_\alpha$$
    a contradiction, then $\aleph_{\alpha+1}$ is a regular cardinal.
\end{proof}

On the other hand, we have
\begin{align*}
\aleph_\omega & = \lim_{n\to\omega}\aleph_n\\
\aleph_{\omega+\omega} & = \lim_{n\to\omega}\aleph_{\omega+n}\\
\aleph_{\omega\cdot\omega} & = \lim_{n\to\omega}\aleph_{\omega\cdot n}\\
\aleph_{\omega_1} & = \lim_{\alpha\to\omega_1}\aleph_\alpha
\end{align*}

It's easy to see that the vast majority of limit cardinals we have contacts with, except for $\aleph_0$, are indeed singular; we will state this more precisely in the next theorem. Furthermore, each cardinal on the examples shown above are singular, the first three had cofinality $\omega$ and the last one cardinality $\omega_1$. We will see how powerful this is when we present other intuitive ways of visualizing them, but for now let us answer some natural questions that arise, such as: are there arbitrarily large singular cardinals? Or rather, are all uncountable limit cardinals singular?

\section{Some Basic Properties}

Let's us first show that there do indeed exist arbitrarily large singular cardinals, and further, if $\aleph_\alpha$ is a singular limit cardinal, then the smallest limit cardinal larger than $\aleph_\alpha$, i.e., the next limit cardinal, is also singular.

\begin{lemma}{}{}
    There are arbitrarily large singular cardinals.
\end{lemma}
\begin{proof}
    Let $\aleph_\alpha$ be an arbitrary cardinal and consider the following sequence
    $$\aleph_\alpha,\aleph_{\alpha+1},\aleph_{\alpha+2},\dots,\aleph_{\alpha+n},\dots$$
    For $n\in\omega$. Therefore
    $$\aleph_{\alpha+\omega}=\lim_{n\to\omega}\aleph_{\alpha+n}$$
    and, then, $\aleph_{\alpha+\omega}$ is a singular cardinal greater than $\aleph_\alpha$.
\end{proof}

As a scholium of the previous lemma, we have

\begin{corollary}{}{}
    If $\aleph_\alpha$ is a singular cardinal, then the next limit cardinal $\aleph_{\alpha+\omega}$ is also singular. In particular, if $\aleph_{\alpha}$ was regular, the next limit cardinal is also singular.
\end{corollary}

With this, we can have a visualization of which of the first cardinals are regular or singular. We have that, after $\aleph_0$, each successor cardinal below $\aleph_\omega$ is regular, and $\aleph_\omega$ is itself singular. We can repeat it from $\aleph_\omega$ to $\aleph_{\omega\cdot 2}$, $\aleph_{\omega^\omega}$, $\aleph_{\varepsilon_0}$, and so on. In fact, even reaching $\aleph_{\omega_1}$, which is singular, we will have the same between it and the next limit $\aleph_{\omega_1+\omega}$. So the ``line'' of cardinals has, apparently, almost always a singular cardinal in the limit points, and always countable many regular successor cardinals between them.

\section{Uncountable Limit Regular Cardinals}\label{weaky_inacces}

Let's now investigate how would be an uncountable limit regular cardinal. Let $\aleph_\lambda$ be such cardinal, since $\lambda$ is a limit ordinal, then
$$\aleph_\lambda=\lim_{\alpha\to\lambda}\aleph_\alpha$$
i.e., $\aleph_\lambda$ is the limit of an increasing sequence of length $\alpha$, in particular, we could change such sequence to another one of length $\text{cf}(\lambda)$. It's easy to see how this imply the next result
\begin{lemma}{}{}
    If $\lambda$ is a limit ordinal, then $\text{cf}(\aleph_\lambda)=\text{cf}(\lambda)$
\end{lemma}

Now, since in our case $\aleph_\lambda$ is also a regular cardinal, then $\aleph_\lambda=\text{cf}(\aleph_\lambda)=\text{cf}(\lambda)\leq\lambda$, which imply
\begin{equation*}
\fbox{$\aleph_\lambda=\lambda$}
\end{equation*}
Such cardinals are called aleph fixed points, since, well, they're a fixed point of the aleph operation.

Now, we already know $\aleph_\omega\neq\omega$, and $\aleph_{\omega_1}\neq\omega_1$, $\aleph_{\omega_2}\neq\omega_2$, etc. and the gap between the indices $\omega_1$ and $1$, $\omega_2$ and $2$, and so on keep increasing! Similarly, we have $\aleph_{\aleph_\omega}\neq\aleph_\omega$. Repeating this arguments it's easy to see that such finite towers of $\aleph$'s cannot reach an aleph fixed point. In fact, it shows that the first aleph fixed point is of the form
$$\aleph_{\aleph_{\aleph_{\small\ddots}}}=\bigcup_{n\in\omega}\aleph_{\aleph_{\aleph_{\small\ddots_{\aleph_0}}}}$$
where we have $n$ $\aleph$'s on the RHS.

Such property and example suggest these cardinals should be really big. However, the next Lemma reveals us that, in fact, this property isn't too strong as it seems, since we can actually prove that there are arbitrarily large singular aleph fixed points, as the \todo{add footnote hyperlink}next lemma\footnote{In fact, this lemma is the toy version of a more general fact known as The Fixed-Point Lemma for normal functions, which guarantees us that every normal function has arbitrarily large fixed points. You can see this in its full generality on my Large Cardinals notes, the proof follows the same idea as the one used to prove the lemma.} states.

\begin{lemma}{}{}
    There are arbitrarily large singular aleph fixed points.
\end{lemma}

\begin{proof}
    Let $\kappa$ be an arbitrary cardinal and consider the following sequence $\alpha_n$ defined recursively as $\alpha_0=\aleph_\gamma$, $\alpha_{n+1}=\aleph_{\alpha_n}$, for each $n\in\omega$. Now, assume by contradiction that there is not singular $\alpha\geq\kappa$ which is a aleph fixed point, then $\alpha_n$ is an increasing sequence; let $\alpha:=\lim_{n\to\omega}\alpha_n$. Obviously the sequence $\aleph_{\alpha_n}=\alpha_{n+1}$ also has limit $\aleph_\alpha$, therefore, we have that
    $$\aleph_\alpha=\lim_{n\to\omega}\aleph_{\alpha_n}=\lim_{n\to\omega}\alpha_{n+1}=\alpha$$
    Then $\alpha=\aleph_\alpha$ is a fixed aleph point which is also a limit of an increasing sequence of length $\omega$, i.e., $\text{cf}(\alpha)=\omega<\alpha$, then it's singular.
\end{proof}

The next natural step is then ask what would be an example of an uncountable regular limit cardinal. Well, surprisingly, such cardinals, named weakly inaccessible cardinals, are so big that ZFC can't prove they exist! Inaccessible cardinals are some of the first cardinals in a larger class of properties named Large Cardinals Properties.

Unfortunatelly, the proof is not that easy, since we need to use a consistent result, namely, that GCH is consistent with ZFC\footnote{You can find a proof of it in my \href{https://github.com/Xennonio/Consistency-Results}{Consistency Results notes}, in Section 3. The Constructible Universe $L$.}. However, we can prove it for a somewhat similar cardinal, named strongly limit cardinals, you can see how to prove it in the \todo{add hyperlink}Large Cardinals notes, and then how to use it to show ZFC also can't prove weakly inaccessible exist.

\section{Rambling on Cofinality}\label{ramb_cof}

Now we already motivated Large Cardinal, let's ramble a little about cofinality. Some books, before defining cofinality, use the following definitions: a set $C\subseteq\lambda$ is cofinal or unbounded in $\lambda$ if for all $\alpha<\lambda$, there exists a $\gamma\in C$ such that $\alpha<\gamma$. It's not hard to see that this is equivalent to $\bigcup C=\lambda$.

We can then define $\cf(\lambda)$ as the least cardinality of an unbounded subset of $\lambda$. In other words, what we are saying is that subsets of regular cardinal need almost all elements to become unbounded in it, whereas singular cardinals don't.

Now, imagine we have some property about ordinals smaller than $\aleph_\omega$, and such property is somewhat inductive. Then the process of extending this property to $\aleph_\omega$ itself would require only $\omega$ many steps, since $\text{cf}(\aleph_\omega)=\omega$, whereas, for example, doing the same for $\aleph_1$ would require uncountable many steps.

Of course none of this is really formal, but you will get your own intuition about this while getting deeper and deeper in set theory. You will see that a lot of theorems states something only for regular or only for singular cardinals, and that they influence a lot, for example, in cardinal arithmetic. In fact, we will see in the next chapters that cardinal arithmetic depends a lot of when the cardinals we are using are either singular or regular.

\section{Additional Properties of Cofinality}

Let's finish this section with some extra results about cofinality.

We already proved that if $\aleph_\alpha$ is singular, then $\cf(\omega_\alpha)<\omega_\alpha$ and, if it's regular, then $\cf(\omega_\alpha)=\omega_\alpha$. We will now show that

\begin{lemma}{}{}
    If a limit ordinal $\alpha$ isn't a cardinal, then $\text{cf}(\alpha)<\alpha$. In particular, for all limit ordinal $\alpha$, we have $\text{cf}(\alpha)=\alpha$ if, and only if, $\alpha$ is a regular cardinal.
\end{lemma}

\begin{proof}
    Let $\alpha$ be an ordinal which is not a cardinal, then $\kappa=\lvert\alpha\rvert<\alpha$ and there is a bijection $f:\kappa\to\alpha$, i.e., a sequence $\langle\alpha_\nu=f(\nu):\nu<\kappa\rangle$ of length $\kappa$ such that $\{\alpha_\nu:\nu<\kappa\}=\alpha$.

    Now, even if such sequence was not increasing, we could find, using transfinite induction, an increasing subsequence with limit $\alpha$. And, since the length of the subsequence is at most $\kappa<\alpha$, then $\text{cf}(\alpha)<\alpha$.
\end{proof}

\begin{lemma}{}{}
    For all limit ordinal $\alpha$, we have
    $$\cf(\cf(\alpha))=\cf(\alpha)$$
\end{lemma}

\begin{proof}
    Let $\vartheta=\text{cf}(\alpha)$. Clearly $\vartheta$ is a limit ordinal and we already know that $\text{cf}(\vartheta)\leq\vartheta$. Now, assume by a contradiction that $\gamma=\text{cf}(\vartheta)<\vartheta$, then there is an increasing sequence $\langle\nu_\xi:\xi<\gamma\rangle$ with $\lim_{\xi\to\gamma}\nu_\xi=\vartheta$. But since $\vartheta=\text{cf}(\alpha)$, there is also a $\langle\alpha_{\nu_\xi}:\xi<\gamma\rangle$ of length $\gamma<\vartheta$ with limit $\alpha$, contradicting the minimality of $\vartheta$.
\end{proof}

\begin{corollary}{}{}
    For all limit ordinal $\alpha$, $\text{cf}(\alpha)$ is a regular cardinal.
\end{corollary}